% Generated draft: requires mathematical and source review.
\documentclass[11pt]{article}
\usepackage[margin=25mm]{geometry}
\usepackage{fontspec}
\setmainfont{Latin Modern Roman}
\usepackage{amsmath,amssymb,mathtools,unicode-math}
\setmathfont{Latin Modern Math}
\usepackage{enumitem,needspace,xurl,hyperref}
\setlength{\parindent}{0pt}
\setlength{\parskip}{5pt}
\setlength{\emergencystretch}{3em}
\newcommand{\N}{\mathbb{N}}
\newcommand{\Z}{\mathbb{Z}}
\newcommand{\Q}{\mathbb{Q}}
\newcommand{\R}{\mathbb{R}}
\newcommand{\C}{\mathbb{C}}
\newcommand{\F}{\mathbb{F}}
\newcommand{\A}{\mathbb{A}}
\newcommand{\PP}{\mathbb{P}}
\newcommand{\E}{\mathbb{E}}
\newcommand{\eps}{\varepsilon}
\newcommand{\dd}{\,\mathrm{d}}
\begin{document}
\section*{UnsolvedMath ID: 20001144}
\textbf{Source problem code:} catalogue entry
\subsection{Definitions and mathematical statement}
Let $r\ge1$, let $z_1,\ldots,z_r$ lie on $\mathbb T=\{z\in\C:|z|=1\}$, and let $p_j\ge0$ satisfy $\sum_jp_j=r$. Omit zero-weight factors, so no $0^0$ convention is needed, and put
\[
M=\max_{z\in\mathbb T}\prod_{j:p_j>0}|z-z_j|^{p_j}.
\]
The conjecture, brought by V. Lev and attributed to Konyagin and the presenter, is $M\ge2$. The original remark says the constant is optimal and records the equal-weight and equally spaced special cases and a polynomial reformulation.

Combine repeated positive-weight locations into distinct active points $\zeta_1,\ldots,\zeta_m$ with aggregate weights $w_i>0$, $\sum_iw_i=r$. Necessarily $1\le m\le r$. The attached aid attributes to Dubinin's sharp hedgehog-capacity inequality, through a published weighted conformal-map correspondence, the stronger result
\[
M\ge2^{r/m}\ge2.
\]
It identifies equality in the first bound with a regular $m$-gon of active locations and equal aggregate weights $r/m$. Consequently equality $M=2$ in the original normalization requires $m=r$, the original points a regular $r$-gon and all weights $1$. This includes the one-point case; for two entries equality occurs exactly for antipodal points of weights $1,1$. Repeated locations or zero weights reduce $m$ and give a strictly larger lower bound when $m<r$.

For the polynomial version let $P$ be monic of degree $d\ge1$ with all zeros on $\mathbb T$ and exactly $m$ distinct zeros. The aid reports
\[
\|P\|_{\mathbb T}:=\max_{|z|=1}|P(z)|\ge2^{d/m}.
\]
The unit-circle-root hypothesis is essential; this is not asserted for arbitrary monic polynomials.

If the active support lies in a closed arc of angular length $0\le L\le2\pi$, the additional reported bound is
\[
M\ge\max\left\{2^{r/m},\,[2\cos(L/4)]^r\right\}.
\]
The second term comes from the point opposite the arc midpoint. The aid supplies a full-solution literature claim for the original inequality, while exact bibliographic statements of the capacity result and correspondence remain absent.
\subsection{Short English statement}
Maximize a weighted product of distances from prescribed circle points. The supplied aid reports a sharp bound in terms of distinct active support, covering zero weights and repetitions, and gives polynomial and arc-concentration consequences.
\subsection{Sources}
\begin{itemize}
\item Local UnsolvedMath registry, permanent ID $20001144$, source code catalogue entry; canonical \texttt{aim-combinatorics-notes.json}, \texttt{notes[268]}.
\item American Institute of Mathematics, \emph{Recent trends in additive combinatorics}, problem 4.13, brought by V. Lev and originally stated by Konyagin and the presenter: \url{https://aimath.org/WWN/additivecomb/additivecomb.pdf}. The original remark records optimality, the two special cases and polynomial reformulation; no complete literature reference is supplied. No PDF inspection is claimed here.
\item Attached research attempt 1: full solution, recorded run and registry status solved, rigor proved, novelty confidence low, classification SOLVED-IN-LITERATURE. It names Dubinin's hedgehog-capacity inequality and a published weighted conformal-map correspondence without exact identifiers. The machine-generated aid is not peer reviewed; independently verify before citation.
\end{itemize}
\end{document}
